\SourceLeaf{15}
\[
-\frac{\mathrm{d}s'.\Phi r'.\cos MQ'}{r'r'}-\frac{\mathrm{d}s''.\Phi r''.\cos MQ''}{r''r''}-\frac{\mathrm{d}s'''.\Phi r'''.\cos MQ'''}{r'''r'''}-\text{ etc.}
\]
\par
in casu posteriori vero per eandem formulam adiecta parte
\[
-\mathrm{d}\Sigma.\Phi0
\]
\par
Multiplicando hanc expressionem per $\cos MX$, habebimus vim, qua partes corporis intra conum sitae attrahunt punctum in directione axi coordinatarum $x$ parallela atque opposita. Hinc vis, qua corpus integrum agit in eadem directione, exprimetur per integrale $-\int\frac{\mathrm{d}s.\Phi r.\cos MQ.\cos MX}{rr}$, per totam corporis superficiem extensum, siquidem punctum attractum iacet extra corpus, sed adiicere adhuc oportet integrale $-\Phi0.\int\mathrm{d}\Sigma.\cos MX$ per totam superficiem sphaericam extensum, quoties $M$ iacet intra corpus. Nullo porro negotio perspicitur, in casu eo, ubi $M$ iaceat in corporis superficie, adiiciendum quidem esse idem integrale $-\Phi0.\int\mathrm{d}\Sigma.\cos MX$, sed per dimidiam tantummodo sphaerae superficiem extensum, et quidem per hemisphaerium id, quod definitur plano corporis superficiem in $M$ tangente atque ab eadem plani parte iacet, a qua est soliditas corporis in puncto $M$. Ut valorem huius integralis determinemus, concipiamus solidum intra hemisphaerium istud atque planum inclusum. Denotet $\theta$ indefinite angulum inter rectam superficiei huius solidi normalem extrorsumque directam atque rectam axi coordinatarum $x$ parallelam. Hinc per Theorema Primum integrale $\int\mathrm{d}s.\cos\theta$ per \textit{totam} solidi superficiem extensum evanescit, unde si integrale per solam partem planam superficiei extensum supponitur $=J$, integrale per superficiei partem curvam debebit esse $=-J$. Sed in parte curva $\mathrm{d}s$ convenit cum nostro $\mathrm{d}\Sigma$, $\theta$ vero fit $=180^0-MX$. Hinc patet, integrale $-\int\mathrm{d}\Sigma.\cos MX$, per hemisphaerium extensum fieri $=-J$. In parte plana autem superficiei manifesto $\theta$ est constans, atque aequalis valori ipsius $QX$ in puncto $M$, unde $J$ aequalis erit producto cosinus huiusce anguli in aream plani, quae fit $=\pi$. Hinc colligitur, integrale $-\Phi0.\int\mathrm{d}\Sigma.\cos MX$, per hemisphaerium quod supra definivimus extensum, fieri $=-\pi\Phi0.\cos QX$, sumto pro $QX$ valore in puncto $M$. Prorsus eodem modo valor integralis $-\Phi0.\int\mathrm{d}\Sigma.\cos MX$ per hemisphaerium alterum extensum invenitur $=+\pi\Phi0.\cos QX$, unde integrale per totam sphaeram fit $=0$. Ex his omnibus colligimus
